% Physical unit 028; accepted Session 03 after fresh 400-dpi authority replay.
\GordanUnitStart{028}{printed page 23}
\GordanTopFolio{23}
\begingroup\small
\caput{VI}
\sect{1}

E formulis capitis II ab ill. Legendre propositis gravissima momenta lineae geodeticae deduci posse, vidimus.

Ad computandas autem longitudines et extensiones ex $l$, longe aptiores sunt formulae ab ill. Jacobi evolutae in libello:
\[
\text{\gq{Solution nouvelle d'un problème de Géodésie}}
\]
transeundo ab integralibus ellipticis ad functiones ellipticas.

\sect{2}

Quem sequentes loco anguli $\varphi$ per se variabilis inducimus variabilem:
\[
 u=\int_0^\varphi {d\varphi\over \sqrt{1-x^2\sin^2\varphi}}
\]
cujus functionibus omnes quantitates exprimamus. Proinde ad computandum $\varphi$ formulis, inversa aequatione integrali nascentibus, utimur; ad reliquas vero computandas quantitates, referri necesse est hos valores ad aequationes capitis II.

Inde fit:
\[
\begin{gathered}
 \varphi=\operatorname{am}u,\qquad \sin\varphi=Su,
 \qquad \cos\varphi=Cu,\qquad \sqrt{1-x^2\sin^2\varphi}=\Delta u,\\
 \sin l=\sin l_0\,Cu,\qquad
 \cos l=\cos l_0\sqrt{1+\tg^2 l_0\,S^2u},\\
 \tg l=\tg l_0{Cu\over \sqrt{1+\tg^2l_0\,S^2u}},\\
 \sqrt{1-e^2\cos^2l}={b\over a}\cdot {\Delta u\over x'}
 =\sqrt{1-e^2\cos^2l_0}\,\Delta(u),\\
 \tg B=\tg B_0{Cu\over \sqrt{1+\tg^2l_0\,S^2u}},\\
 \sin B=\sin B_0{Cu\over \Delta u},
 \qquad
 \cos B=\cos B_0{\sqrt{1+\tg^2l_0\,S^2u}\over \Delta u},\\
 \sqrt{1-e^2\sin^2B}={\cos B\over \cos l}={x'\over \Delta u}.
\end{gathered}
\]
\endgroup
